\documentclass[11pt]{article}
\usepackage[margin=1in]{geometry}
\usepackage{amsmath,amssymb,amsthm}
\usepackage{booktabs}
\usepackage{listings}
\usepackage{xcolor}
\usepackage[numbers]{natbib}
\usepackage{hyperref}
\sloppy

\newtheorem{theorem}{Theorem}
\newtheorem{lemma}[theorem]{Lemma}
\theoremstyle{definition}
\newtheorem{proposal}[theorem]{Proposed statement}
\theoremstyle{remark}
\newtheorem{remark}[theorem]{Remark}

\lstdefinelanguage{lean}{
  morekeywords={theorem,lemma,suffices,constructor,def,example,import,namespace,end,open,by,intro,exact,obtain,refine,have,simp,only,at,fun,noncomputable},
  sensitive=true, morecomment=[l]{--}, morecomment=[s]{/-}{-/}, morestring=[b]"}
\lstset{language=lean, basicstyle=\ttfamily\footnotesize, keywordstyle=\color{blue!60!black}\bfseries,
  commentstyle=\color{green!40!black}, columns=fullflexible, breaklines=true, frame=single,
  extendedchars=true, inputencoding=utf8,
  literate={≤}{{$\le$}}1 {≥}{{$\ge$}}1 {→}{{$\to$}}1 {↔}{{$\leftrightarrow$}}1 {ℝ}{{$\mathbb{R}$}}1
           {ℕ}{{$\mathbb{N}$}}1 {⟨}{{$\langle$}}1 {⟩}{{$\rangle$}}1 {∃}{{$\exists$}}1 {∀}{{$\forall$}}1
           {∧}{{$\wedge$}}1 {·}{{$\cdot$}}1 {−}{{$-$}}1 {Δ}{{$\Delta$}}1 {μ}{{$\mu$}}1 {σ}{{$\sigma$}}1
           {‖}{{$\|$}}1 {∑}{{$\sum$}}1 {≠}{{$\ne$}}1 {∈}{{$\in$}}1 {←}{{$\leftarrow$}}1 {₀}{{$_0$}}1 {η}{{$\eta$}}1}

\title{Re-certifying the Tempo Cadence Kernel of Viridis OS:\\ A Lean Proof of its Six-Part Cost-Model Core}
\author{Justin Hart\\ Viridis LLC}
\date{21 November 2026}

\begin{document}
\maketitle

\begin{abstract}
This note re-certifies the \emph{Tempo} kernel of Viridis OS, which recommends a monitoring cadence (the interval between
stewardship visits). The result certifies the mathematics the kernel applies and not any field input or stewardship outcome.
Given a haste cost $h>0$ and a neglect cost $n>0$, the kernel assumes that the cost per unit time of visiting every $t>0$
time units is $\Sigma(t)=h/t+nt$. The cost model and its two coefficients are inputs that a user declares; nothing here
shows that they describe any stewardship programme. Until now, the kernel's formal assurance rested on a single receipt
from the retired Comparator verifier, for a Lean module \texttt{StewardshipTempo} whose source is not on file in the canon
repository. A sealed statement, \texttt{P2.Tempo}, was therefore reconstructed from the kernel's reference runner and
from the theorem signatures of a deposited file as relayed by a web tool. We give a Lean proof of that statement under
Lean v4.28.0 and a pinned Mathlib. It has six parts, one for each theorem named in the kernel's registry entry:
(1) $\Sigma$ is strictly convex on $(0,\infty)$;
(2) $\Sigma(t)\ge 2\sqrt{hn}$ for every $t>0$;
(3) the floor is attained at $t^\star=\sqrt{h/n}$;
(4) $t^\star$ is the only positive cadence that attains it;
(5) $(u+1/u)/2=\cosh\log u$ for $u>0$, and so the normalised cost $\Sigma(t)/(2\sqrt{hn})$ equals the haste--neglect tax $\cosh\log(t/t^\star)$;
(6) the normalised cost also equals $\cosh\log(t^\star/t)$, and the tax is at least $1$ and equals $1$ exactly at $t^\star$.
The proof substitutes $h=s^2$, $n=r^2$ and rests on the gap identity $\Sigma(t)-2sr=(s-rt)^2/t$. Strict convexity
is proved directly from the definition, via $(ah/x+bh/y)(ax+by)=h+hab(x-y)^2/(xy)$ for $a+b=1$.
The proof uses only the axioms \texttt{propext}, \texttt{Classical.choice} and \texttt{Quot.sound}; no automated prover was used.
Together the six parts settle that the runner's optimum, floor, normalised cost and tax are mutually consistent consequences
of the declared model, and that over- and under-shooting the optimum by the same factor cost the same.
They do not settle the form of the model, the values or units of $h$ and $n$, the runner's floating-point tolerance,
or whether the deployed code matches the runner. This note is part of a program that brings each live decision kernel
under sealed, re-checkable statements; Lean proofs for four other kernels (Carbon Continuity, Restoration, Afforestation and
Harmonization) from earlier sessions are retained.
\end{abstract}

\section{Introduction}
Viridis OS runs a set of decision kernels. Each one is a small piece of software that turns ecological
parameters into a decision. A kernel's rule is only as trustworthy as the mathematics behind it.
The archived canon~\cite{HartCanon2026} proved those rules in Lean, but the certificates came from tooling that has
since been retired. The project described here moves each kernel's assurance onto a re-checkable footing in two steps.
First, a \emph{sealed statement} fixes the exact proposition a kernel relies on.
Second, a Lean proof of that statement is checked under Lean v4.28.0 with a pinned Mathlib~\cite{Mathlib2020}.

The two steps are separate on purpose. Writing the statement is where most of the risk lies.
A statement can be correct and easy to prove and still be the wrong one. It can be weaker than what the kernel uses,
stronger than what the canon proved, or vacuous. The 30 October revision covered that first step for three kernels:
Restoration, Afforestation and Harmonization.
The 31 October revision took the second step for Restoration (Section~\ref{sec:restoration}), the 1 November revision
took it for Afforestation (Section~\ref{sec:afforestation}), and the present revision takes it for Harmonization
(Section~\ref{sec:harmonization}).
The Restoration kernel makes a go/no-go call for a planted cluster. Go means the cluster is past the nucleation barrier
and will grow into forest. No-go means it is still on the uphill side and is expected to dissolve.
A wrong rule here would send planting budgets to clusters that cannot persist, so the geometry of the barrier is worth
checking mechanically, however elementary it is.
The Afforestation kernel works one step later: given that clusters must clear the barrier, it compares seeding plans and
site-preparation levers (such as nurse plants or mulch, which lower the edge hostility $\sigma$) by how much established canopy each seed buys,
and it scores a proposed seed allocation by its alignment with the optimal one.
Both rules rest on two short facts, a scaling law and a Cauchy--Schwarz bound, which are what \texttt{P2.Afforestation} seals.
The Harmonization kernel addresses coordination: rather than issuing $N$ separate commands, it broadcasts one number,
a price on effort, and lets each steward respond locally. This is only sound if that price both meets the shared target and
leads to an allocation no worse than a central planner's; \texttt{P2.Harmonization} seals exactly those two properties.
The Tempo kernel (Section~\ref{sec:tempo}) answers a scheduling question: how often should a steward visit or intervene?
Acting too often has a cost (haste), and so does waiting too long (neglect).
Unlike the four kernels above, it had no Lean source in the canon repository, so its statement had to be reconstructed from
what the kernel computes (20 November revision); the present revision gives a Lean proof of the sealed statement (Section~\ref{sec:tempoproof}).
The Tempo kernel's earlier assurance rested on a receipt (identifier \texttt{134d3df7-25b5-4ab5-80e5-f747db6e903f}) from the
retired Comparator verifier, for a module whose source is not on file in the canon repository; such a receipt cannot be re-checked.
The question is narrow. We ask what each kernel's rule \emph{mathematically} relies on, and how to write that down
self-contained and faithful to the canon theorem it replaces. We do not ask whether the models describe any real landscape.

\section{Background and related work}
\paragraph{Carbon Continuity.} One disturbance-recovery cycle maps $(L,D)\mapsto(aL+rD,\ pL+dD)$.
This is the nonnegative matrix $M=\begin{pmatrix}a&r\\p&d\end{pmatrix}$, and the threshold theorem is a
two-dimensional instance of the Collatz--Wielandt characterisation in Perron--Frobenius theory~\cite{BermanPlemmons1994}.

\paragraph{Restoration and Afforestation.} The canon models a planted cluster as a nucleus of the forest phase inside a
metastable degraded phase. It uses the free-energy form of classical nucleation theory~\cite{Kashchiev2000}:
a bulk gain $-\Delta\mu\, n$ competes with a surface cost $a\sigma\, n^{2/3}$.
In the cube-root variable $r=n^{1/3}$ this becomes the cubic $G(r)=-\Delta\mu r^3+a\sigma r^2$.
The textbook facts are a critical size at the barrier top and a barrier equal to half the bulk term at the critical size.
These are the mathematical content of \texttt{P2.Restoration}.
The Afforestation efficiency $e=C/n^\star$ inherits the cube from $n^\star=(r^\star)^3$.
The alignment score is a squared cosine, bounded by Cauchy--Schwarz.

\paragraph{Harmonization.} The scalar-price theorem is the separable quadratic case of Lagrangian duality with
one linear coupling constraint~\cite{BoydVandenberghe2004}.
At the clearing price, the decentralised best responses solve the centralised problem, and there is no duality gap.
In economics this is the first welfare theorem in its simplest setting, a single resource constraint allocated among
agents with separable quadratic costs; in engineering it is the scalar-price case of dual decomposition.

\paragraph{Tempo.} The cost $h/t+nt$ trades a term that falls with the interval against one that rises with it.
Its minimisation is the two-term arithmetic--geometric mean inequality, and $h/t$ is a standard example of a strictly convex
function on $(0,\infty)$~\cite{BoydVandenberghe2004}.
The canon's working paper on the kernel is the Zenodo record~\cite{HartTempo2026}.

We claim no novelty for any of this mathematics. The contribution is a set of faithful, machine-checkable targets tied
to deployed kernels.

\section{Carbon Continuity (P2-T001)}
The sealed module defines $\mathrm{stepLiving}(a,r,L,D)=aL+rD$ and $\mathrm{stepDurable}(d,p,L,D)=pL+dD$.
$\mathrm{Sustainable}(a,d,r,p)$ means that some $L,D>0$ satisfy $L\le aL+rD$ and $D\le pL+dD$.
\begin{theorem}[\texttt{P2.CarbonContinuity}]
For all real $a,d,r,p$ with $0\le a<1$, $0\le d<1$, $r>0$ and $p>0$:
$\mathrm{Sustainable}(a,d,r,p)\iff (1-a)(1-d)\le rp$.
\end{theorem}
\begin{proof}[Proof sketch]
Sufficiency uses the witness $(L,D)=(r,1-a)$. The living step is then an identity, and the durable step reduces to
$rp-(1-a)(1-d)\ge 0$. For necessity, rearrange the two inequalities to $(1-a)L\le rD$ and $(1-d)D\le pL$,
multiply them (both sides are positive), and cancel $LD>0$.
\end{proof}
The previous revision gave a Lean proof of this statement (file \texttt{proofs/P2.CarbonContinuity.lean} of run
2026-10-29), with axioms \texttt{propext}, \texttt{Classical.choice} and \texttt{Quot.sound}.
\begin{lstlisting}
theorem main : ViridisStatements.P2.CarbonContinuity.Statement := by
  intro a d r p _ ha1 _ hd1 hr _
  exact ⟨necessary a d r p ha1 hd1, sufficient a d r p ha1 hr⟩
\end{lstlisting}
We observe that the helper lemmas use only $a<1$, $d<1$ and $r>0$.

\section{Restoration: the critical nucleus (P2-T003)}\label{sec:restoration}
\subsection{The sealed statement}
The sealed module \texttt{ViridisStatements.P2.Restoration} (sha256 \texttt{fd40b113\ldots ccff18}) is the statement
proposed on 30 October (Section~\ref{sec:proposals}), unchanged.
It defines $r^\star(\Delta\mu,a,\sigma)=2a\sigma/(3\Delta\mu)$ and $G(\Delta\mu,a,\sigma,r)=-\Delta\mu r^3+a\sigma r^2$
verbatim from the canon file \texttt{ForestNucleation.lean} (sha256 \texttt{667a826a\ldots ec28132}, which we checked
against the pinned manifest).
\begin{theorem}[\texttt{P2.Restoration}]\label{thm:restoration}
\begin{enumerate}
\item[(i)] For all real $\Delta\mu,a,\sigma>0$: $r^\star>0$, $G(r)\le G(r^\star)$ for every $r\ge 0$, and every $r\ge0$ with
$G(r)=G(r^\star)$ equals $r^\star$.
\item[(ii)] For all real $\Delta\mu>0$ and all real $a,\sigma$: $G(r^\star)=\tfrac12\Delta\mu\,(r^\star)^3$.
\end{enumerate}
\end{theorem}
We give a Lean proof of Theorem~\ref{thm:restoration} in \texttt{proofs/P2.Restoration.lean}.
The local checker, which runs the same stages as the verifier, passes it, and the only axioms it reports are
\texttt{propext}, \texttt{Classical.choice} and \texttt{Quot.sound}.

\subsection{Proof}
\begin{proof}
Let $\Delta\mu\ne0$. By the definition of $r^\star$, $a\sigma=\tfrac32\Delta\mu\, r^\star$.
Substituting into $G$ gives $G(r)=\Delta\mu\bigl(-r^3+\tfrac32 r^\star r^2\bigr)$, hence
$G(r^\star)=\tfrac12\Delta\mu (r^\star)^3$, which is (ii). Subtracting,
\[
G(r^\star)-G(r)=\Delta\mu\bigl(r^3-\tfrac32 r^\star r^2+\tfrac12 (r^\star)^3\bigr)=\Delta\mu\,(r-r^\star)^2\,(r+r^\star/2).
\]
Now let $\Delta\mu,a,\sigma>0$. Then $r^\star>0$ as a quotient of positive numbers.
For $r\ge0$ all three factors on the right are nonnegative, so $G(r)\le G(r^\star)$.
If equality holds, the product is zero. Since $\Delta\mu>0$ and $r+r^\star/2>0$, we get $(r-r^\star)^2=0$, that is, $r=r^\star$.
\end{proof}
The Lean development follows this argument line by line. Its key lemmas are:
\begin{lstlisting}
lemma aσ_eq (Δμ a σ : ℝ) (hΔμ : Δμ ≠ 0) :
    a * σ = 3 / 2 * Δμ * rStar Δμ a σ := by
  unfold rStar
  field_simp

lemma key_identity (Δμ a σ r : ℝ) (hΔμ : Δμ ≠ 0) :
    G Δμ a σ (rStar Δμ a σ) - G Δμ a σ r =
      Δμ * (r - rStar Δμ a σ) ^ 2 * (r + rStar Δμ a σ / 2) := by
  unfold G
  rw [aσ_eq Δμ a σ hΔμ]
  ring

lemma barrier_half (Δμ a σ : ℝ) (hΔμ : 0 < Δμ) :
    G Δμ a σ (rStar Δμ a σ) = (1 / 2) * Δμ * (rStar Δμ a σ) ^ 3 := by
  unfold G
  rw [aσ_eq Δμ a σ hΔμ.ne']
  ring

theorem main : ViridisStatements.P2.Restoration.Statement :=
  ⟨fun Δμ a σ hΔμ ha hσ => critical_nucleus Δμ a σ hΔμ ha hσ,
   fun Δμ a σ hΔμ => barrier_half Δμ a σ hΔμ⟩
\end{lstlisting}
Here \texttt{critical\_nucleus} proves (i). Maximality comes from \texttt{key\_identity}, \texttt{positivity} and \texttt{linarith}.
Uniqueness is a two-way \texttt{mul\_eq\_zero} case split followed by \texttt{pow\_eq\_zero\_iff}.

\subsection{Relation to the canon proof}
The canon proof uses the same factorisation, but in a compressed tactic style (\texttt{refine'}, a long term-mode
\texttt{contrapose!} chain, and \texttt{grind} for the half identity).
We observe that the canon proofs, pasted verbatim into a candidate file, also pass the local checker under the pinned
toolchain. The archived proof therefore still replays.
We submit the restructured version because its helper lemmas name the single algebraic fact everything rests on,
and because it depends less on the behaviour of \texttt{grind}, which may change between Mathlib versions.

\subsection{Observations}
\begin{itemize}
\item \textbf{An unused hypothesis.} Part (ii) holds for every real $\Delta\mu$. For $\Delta\mu\ne0$ the argument above
does not use the sign. For $\Delta\mu=0$, Lean's convention $x/0=0$ gives $r^\star=0$, and both sides are $0$.
We checked both cases locally. The sealed statement keeps $0<\Delta\mu$ to match the canon.
\item \textbf{The sign of $\Delta\mu$ matters for (i).} If $\Delta\mu<0$ the degraded state is favoured, $r^\star<0$, and
$G$ is unbounded above on $r\ge0$, so no cluster size is a barrier top. The sealed module includes a negative control for this case.
\item \textbf{$r$ versus $n$.} The theorem concerns the cube-root variable $r=n^{1/3}$. Since $r\mapsto r^3$ is an increasing
bijection of $[0,\infty)$, it follows informally that $n^\star=(r^\star)^3$ is the unique maximiser of
$-\Delta\mu n+a\sigma n^{2/3}$ on $n\ge0$. This translation is not part of the sealed statement, and it has no Lean proof here.
\end{itemize}

\section{Afforestation: the seeding core (P2-T004)}\label{sec:afforestation}
\subsection{The sealed statement}
The sealed module \texttt{ViridisStatements.P2.Afforestation} (sha256 \texttt{2fc1cbb0\ldots 59436}) is byte-identical to the
statement proposed on 30 October (Section~\ref{sec:proposals}).
It defines $n^\star(a,\sigma,\Delta\mu)=(2a\sigma/(3\Delta\mu))^3$, $\mathrm{eff}(a,\sigma,\Delta\mu,C)=C\,(3\Delta\mu/(2a\sigma))^3$
and, on any real inner-product space $E$, $\eta(n,m)=\langle n,m\rangle^2/(\|n\|^2\|m\|^2)$.
The first two are verbatim from the canon file \texttt{AfforestationStewardship.lean} (sha256 \texttt{3ebe33bf\ldots 9526f},
checked against the pinned manifest); $\eta$ names the expression in the canon theorem's conclusion.
\begin{theorem}[\texttt{P2.Afforestation}]\label{thm:afforestation}
\begin{enumerate}
\item[(i)] For all real $a,\sigma,\Delta\mu,t>0$ and every real $C$:
$\mathrm{eff}(a,\sigma,t\Delta\mu,C)=t^3\,\mathrm{eff}(a,\sigma,\Delta\mu,C)$.
\item[(ii)] For every real inner-product space $E$ and all $n,m\in E$ with $n\ne0$, $m\ne0$: $0\le\eta(n,m)\le1$.
\end{enumerate}
\end{theorem}
We give a Lean proof of Theorem~\ref{thm:afforestation} in \texttt{proofs/P2.Afforestation.lean}.
The local checker passes it, and the only axioms it reports are \texttt{propext}, \texttt{Classical.choice} and \texttt{Quot.sound}.

\subsection{Proof}
\begin{proof}
(i) $C\,(3t\Delta\mu/(2a\sigma))^3=t^3\cdot C\,(3\Delta\mu/(2a\sigma))^3$ is an identity of rational expressions in which
$t$ appears only in the numerator.
(ii) $\eta$ is a square divided by a square, so $\eta\ge0$. By Cauchy--Schwarz, $|\langle n,m\rangle|\le\|n\|\,\|m\|$;
both sides are nonnegative, so squaring gives $\langle n,m\rangle^2\le\|n\|^2\|m\|^2$, and dividing by the nonnegative
denominator gives $\eta\le1$.
\end{proof}
In Lean, division by zero is defined to be zero, so both steps go through without the hypotheses: for (i), \texttt{ring}
treats $x/y$ as $x\cdot y^{-1}$; for (ii), \texttt{div\_le\_one\_of\_le\(_0\)} needs only a nonnegative denominator.
The submitted file therefore proves hypothesis-free helper lemmas and discards the hypotheses in \texttt{main}:
\begin{lstlisting}
theorem eff_cubic (a σ Δμ C t : ℝ) :
    eff a σ (t * Δμ) C = t ^ 3 * eff a σ Δμ C := by
  unfold eff
  ring

theorem eta_mem_unit {E : Type*} [NormedAddCommGroup E] [InnerProductSpace ℝ E]
    (n m : E) : 0 ≤ eta n m ∧ eta n m ≤ 1 := by
  unfold eta
  refine ⟨by positivity, ?_⟩
  apply div_le_one_of_le₀ _ (by positivity)
  have h := abs_real_inner_le_norm n m
  have h2 : |inner ℝ n m| ^ 2 ≤ (‖n‖ * ‖m‖) ^ 2 :=
    pow_le_pow_left₀ (abs_nonneg _) h 2
  rw [sq_abs, mul_pow] at h2
  exact h2

theorem main : ViridisStatements.P2.Afforestation.Statement := by
  refine ⟨fun a σ Δμ C t _ _ _ _ => eff_cubic a σ Δμ C t, ?_⟩
  intro E _ _ n m _ _
  exact eta_mem_unit n m
\end{lstlisting}

\subsection{Relation to the canon proof}
Part (i) is the canon proof unchanged (\texttt{unfold eff; ring}).
For part (ii) the canon closes the numerator bound with \texttt{nlinarith} fed the two sides of
$|\langle n,m\rangle|\le\|n\|\|m\|$; we replace that search by an explicit squaring step
(\texttt{pow\_le\_pow\_left\(_0\)}, \texttt{sq\_abs}, \texttt{mul\_pow}), which does not depend on \texttt{nlinarith}'s heuristics.
We observe that the canon proofs, pasted verbatim (including \texttt{refine'}) into a candidate file, also pass the
local checker under the pinned toolchain.
The sealed statement quantifies over a universe-polymorphic type \texttt{E : Type*}, as the canon does; the polymorphic
candidate passed the statement-match stage.

\subsection{Observations}
\begin{itemize}
\item \textbf{No hypothesis is used.} The cubic law holds for all real $a,\sigma,\Delta\mu,t,C$ (as a Lean identity; when
$a\sigma=0$ both sides are $0$). The alignment bound holds for $n=0$ or $m=0$, where $\eta=0/0=0$ in Lean.
The sealed statement keeps all hypotheses to match the canon. The helper lemmas in the submitted file state the
hypothesis-free versions; they are not sealed statements.
\item \textbf{The name over-promises.} The canon theorem is called \texttt{seeding\_efficiency\_eq\_cos2\_theta}, and its
docstring describes $\eta=\cos^2\Theta$ with forcing debt $1-\eta=\sin^2\Theta$ and $\eta=1$ exactly for aligned allocations.
Only the bound $0\le\eta\le1$ is proved, in the canon and here. A kernel that reports ``zero forcing debt iff the allocation
is optimal'' relies on the equality case of Cauchy--Schwarz, which is not sealed.
\item \textbf{Scale invariance, not proved.} $\eta(cn,m)=\eta(n,m)$ for $c\ne0$, so the alignment score ignores the total
seed budget and measures only the direction of the allocation. This is immediate but has no Lean proof here.
\item \textbf{The site-prep lever.} The canon's 8$\times$ lever (halving $\sigma$ multiplies $\mathrm{eff}$ by $8$) is the same
kind of ring identity; it is pinned in the sealed module only by a concrete \texttt{example}, not as a general statement.
\end{itemize}

\section{Harmonization: the scalar-price welfare core (P2-T005)}\label{sec:harmonization}
\subsection{The sealed statement}
The sealed module \texttt{ViridisStatements.P2.Harmonization} (sha256 \texttt{48934ecb\ldots 4b3dd7}) is byte-identical to the
statement proposed on 30 October (Section~\ref{sec:proposals}).
Its definitions are verbatim from the canon file \texttt{GaianHarmonization.lean} (sha256 \texttt{cfdead65\ldots 110aec},
checked against the pinned manifest): for $a,b\in\mathbb{R}^N$, a price $\lambda$ and a target $X$,
\[
\mathrm{br}(\lambda)_i=b_i-\frac{\lambda}{a_i},\qquad J(x)=\sum_i\frac{a_i}{2}(x_i-b_i)^2,\qquad
\lambda^\star=\frac{\sum_i b_i-X}{\sum_i 1/a_i},\qquad x^\star_i=\mathrm{br}(\lambda^\star)_i .
\]
Here $\mathrm{br}(\lambda)_i$ minimises agent $i$'s private Lagrangian $\tfrac{a_i}{2}(x-b_i)^2+\lambda x$.
\begin{theorem}[\texttt{P2.Harmonization}]\label{thm:harmonization}
For every $N>0$, all $a,b\in\mathbb{R}^N$ with every $a_i>0$, and every real $X$:
\begin{enumerate}
\item[(i)] (market clearing) $\sum_i\mathrm{br}(\lambda^\star)_i=X$;
\item[(ii)] (no duality gap) every $y\in\mathbb{R}^N$ with $\sum_i y_i=X$ satisfies $J(x^\star)\le J(y)$.
\end{enumerate}
\end{theorem}
We give a Lean proof of Theorem~\ref{thm:harmonization} in \texttt{proofs/P2.Harmonization.lean}.
The local checker passes it, and the only axioms it reports are \texttt{propext}, \texttt{Classical.choice} and \texttt{Quot.sound}.

\subsection{Proof}
\begin{proof}
Write $S=\sum_i 1/a_i$; since $N>0$ and every $a_i>0$, $S>0$.
(i) The aggregate best response is affine in the price: $\sum_i\mathrm{br}(\lambda)_i=\sum_i b_i-\lambda S$.
At $\lambda=\lambda^\star$ the second term is $\sum_i b_i-X$, so the sum is $X$.
(ii) Since $a_i\ne0$, $a_i(x^\star_i-b_i)=-\lambda^\star$ for every $i$ (the equimarginal identity).
For each $i$, expanding the square around $x^\star_i$,
\[
\tfrac{a_i}{2}(y_i-b_i)^2=\tfrac{a_i}{2}(x^\star_i-b_i)^2+a_i(x^\star_i-b_i)(y_i-x^\star_i)+\tfrac{a_i}{2}(y_i-x^\star_i)^2
\ \ge\ \tfrac{a_i}{2}(x^\star_i-b_i)^2-\lambda^\star(y_i-x^\star_i),
\]
using $a_i\ge0$. Summing over $i$, the price terms add up to $-\lambda^\star\bigl(\sum_i y_i-\sum_i x^\star_i\bigr)=-\lambda^\star(X-X)=0$
by feasibility of $y$ and part (i). Hence $J(y)\ge J(x^\star)$.
\end{proof}
The Lean development follows this argument. Its key lemmas are:
\begin{lstlisting}
theorem sum_br {N : ℕ} (a b : Fin N → ℝ) (lam : ℝ) :
    (∑ i, br a b lam i) = (∑ i, b i) - lam * ∑ i, 1 / a i := by
  unfold br
  rw [Finset.sum_sub_distrib, Finset.mul_sum]
  congr 1
  refine Finset.sum_congr rfl fun i _ => ?_
  rw [mul_one_div]

theorem market_clears {N : ℕ} (a b : Fin N → ℝ) (hN : 0 < N) (ha : ∀ i, 0 < a i) (X : ℝ) :
    (∑ i, br a b (lamStar a b X) i) = X := by
  rw [sum_br]
  unfold lamStar
  rw [div_mul_cancel₀ _ (sum_inv_pos a hN ha).ne']
  ring

theorem equimarginal {N : ℕ} (a b : Fin N → ℝ) (ha : ∀ i, 0 < a i) (X : ℝ) (i : Fin N) :
    a i * (xStar a b X i - b i) = - lamStar a b X := by
  unfold xStar br
  field_simp [(ha i).ne']
  ring

theorem pointwise {N : ℕ} (a b : Fin N → ℝ) (ha : ∀ i, 0 < a i) (X : ℝ) (y : Fin N → ℝ)
    (i : Fin N) :
    (a i / 2) * (xStar a b X i - b i) ^ 2 + (- lamStar a b X) * (y i - xStar a b X i)
      ≤ (a i / 2) * (y i - b i) ^ 2 := by
  rw [← equimarginal a b ha X i]
  nlinarith [mul_nonneg (ha i).le (sq_nonneg (y i - xStar a b X i))]

theorem strong_duality {N : ℕ} (a b : Fin N → ℝ) (hN : 0 < N) (ha : ∀ i, 0 < a i) (X : ℝ)
    (y : Fin N → ℝ) (hy : (∑ i, y i) = X) :
    J a b (xStar a b X) ≤ J a b y := by
  have hsum := Finset.sum_le_sum fun i (_ : i ∈ Finset.univ) => pointwise a b ha X y i
  have hx : (∑ i, xStar a b X i) = X := market_clears a b hN ha X
  have hzero : (∑ i, (- lamStar a b X) * (y i - xStar a b X i)) = 0 := by
    rw [← Finset.mul_sum, Finset.sum_sub_distrib, hy, hx, sub_self, mul_zero]
  rw [Finset.sum_add_distrib, hzero, add_zero] at hsum
  exact hsum

theorem main : ViridisStatements.P2.Harmonization.Statement := by
  intro N a b hN ha X
  exact ⟨market_clears a b hN ha X, fun y hy => strong_duality a b hN ha X y hy⟩
\end{lstlisting}
Here \texttt{sum\_inv\_pos} proves $S>0$ with \texttt{Finset.sum\_pos}, using the index $0$, which exists because $N>0$.

\subsection{Relation to the canon proof}
The canon proof has the same structure (equimarginal identity, pointwise inequality, summation), but closes the clearing step
with two \texttt{simp +decide} calls over a list of commutativity lemmas, and the final summation with \texttt{convert \ldots\ using 1}
followed by \texttt{simp +decide} and a disjunction.
We replace these by the explicit linearity lemma \texttt{sum\_br} and an explicit cancellation of the price terms, so that the
argument depends less on the simp set, which may change between Mathlib versions.
We observe that the three canon theorems (\texttt{ght\_market\_clears}, \texttt{ght\_equimarginal},
\texttt{ght\_strong\_duality\_decentralized\_eq\_centralized}), pasted verbatim against the sealed definitions and combined in a
\texttt{main}, also pass the local checker under the pinned toolchain.

\subsection{Observations}
\begin{itemize}
\item \textbf{$N>0$ is needed for (i) only.} The sealed module's negative control shows that (i) fails for $N=0$.
Part (ii) holds for $N=0$ as well: both dissipations are empty sums. We checked this locally in Lean; it is not a sealed statement.
\item \textbf{Positive curvature matters for (ii).} Under Lean's convention $1/0=0$, allowing a zero curvature breaks optimality.
With $N=2$, $a=(0,1)$, $b=(0,0)$, $X=1$ one gets $\lambda^\star=-1$, $x^\star=(0,1)$ and $J(x^\star)=\tfrac12$, while the feasible
$y=(1,0)$ has $J(y)=0$. We checked this with \texttt{norm\_num}. The proof uses $a_i\ge0$ in the pointwise inequality and $a_i\ne0$ in the
equimarginal identity, so the hypothesis $a_i>0$ is exactly what is needed.
\item \textbf{What is not sealed.} Uniqueness of the minimiser (true, by strict convexity, since the slack $\tfrac{a_i}{2}(y_i-x^\star_i)^2$ is
positive unless $y=x^\star$), the equimarginal identity as a standalone statement, and the remaining canon results
(the trivial scalar-sufficiency existence statement, the information-bottleneck bandwidth results, the robustness ratios, and the
deferred primal--dual convergence of price adjustment) have no Lean proof against a sealed statement here.
\end{itemize}

\section{Method for writing a sealed statement}
For each kernel we followed the same procedure.
\begin{enumerate}
\item \textbf{Identify the canon theorems.} We take the theorems named in the task brief and check the canon file's SHA-256
against the pinned manifest (canon commit \texttt{80d32904}).
\item \textbf{Copy definitions verbatim.} We copy only the definitions the theorems use, with the same argument order,
and drop the rest. Any added abbreviation must unfold to the canon expression exactly.
\item \textbf{Copy hypotheses verbatim.} Each conjunct of \texttt{Statement} keeps exactly the hypotheses of its canon theorem.
We drop no hypothesis, even one that is mathematically unnecessary, because dropping it would strengthen the canon statement.
We add none, because adding one would weaken it.
\item \textbf{State the conclusion as proved, not as named.} Where a canon docstring or name claims more than the Lean
conclusion, the statement follows the conclusion.
\item \textbf{Pin the meaning.} We add sanity \texttt{example}s with concrete numbers, closed by \texttt{norm\_num} or \texttt{simp}
(\texttt{decide} does not apply to $\mathbb{R}$). They include negative controls showing that the key hypotheses or factors cannot be dropped.
\item \textbf{Test non-triviality.} We compile each module and probe \texttt{Statement} with 28 one-shot tactics
after \texttt{unfold Statement}. The tactics include \texttt{rfl}, \texttt{decide}, \texttt{simp}, \texttt{norm\_num}, \texttt{aesop},
\texttt{linarith}, \texttt{nlinarith}, \texttt{positivity}, \texttt{grind}, \texttt{field\_simp}, \texttt{ring} and \texttt{polyrith}.
\item \textbf{Test feasibility.} Privately, we check that a port of the canon proof closes the proposed statement.
This guards against proposing a false target. It is not a claim of proof.
\end{enumerate}

\section{The three proposals (P2-T002)}\label{sec:proposals}
The three statements below were proposed on 30 October. \texttt{P2.Restoration} has since been sealed unchanged and is proved in
Section~\ref{sec:restoration}; \texttt{P2.Afforestation} has been sealed unchanged and is proved in
Section~\ref{sec:afforestation}; \texttt{P2.Harmonization} has been sealed unchanged and is proved in
Section~\ref{sec:harmonization}.

\subsection{\texttt{P2.Restoration}: the critical-nucleus go/no-go}
Definitions: $r^\star(\Delta\mu,a,\sigma)=2a\sigma/(3\Delta\mu)$ and $G(\Delta\mu,a,\sigma,r)=-\Delta\mu r^3+a\sigma r^2$.
\begin{proposal}
(i) For all $\Delta\mu,a,\sigma>0$: $r^\star>0$; $G(r)\le G(r^\star)$ for all $r\ge0$; and $G(r)=G(r^\star)$ with $r\ge0$ implies $r=r^\star$.
(ii) For all $\Delta\mu>0$ and all $a,\sigma$: $G(r^\star)=\tfrac12\Delta\mu\,(r^\star)^3$.
\end{proposal}
Part (i) restates canon \texttt{critical\_nucleus}, and part (ii) restates \texttt{barrier\_half\_identity}
(file sha256 \texttt{667a826a\ldots ec28132}).
Both follow from the factorisation $G(r^\star)-G(r)=\Delta\mu\,(r-r^\star)^2\,(r+r^\star/2)$.
The kernel's go/no-go rule compares cluster size with $n^\star=(r^\star)^3$. Below $n^\star$ the cluster is on the uphill side of the barrier.
Negative controls: $G(-1)=5/2>G(r^\star)=1/2$ at $(\Delta\mu,a,\sigma)=(1,1,3/2)$, so the restriction $r\ge0$ is essential.
$G(r^\star)\ne\Delta\mu(r^\star)^3$, so the factor $\tfrac12$ is substantive.
$\Delta\mu=-1$ gives $r^\star<0$.
\begin{lstlisting}
def Statement : Prop :=
  (∀ Δμ a σ : ℝ, 0 < Δμ → 0 < a → 0 < σ →
    0 < rStar Δμ a σ ∧
    (∀ r : ℝ, 0 ≤ r → G Δμ a σ r ≤ G Δμ a σ (rStar Δμ a σ)) ∧
    (∀ r : ℝ, 0 ≤ r → G Δμ a σ r = G Δμ a σ (rStar Δμ a σ) → r = rStar Δμ a σ)) ∧
  (∀ Δμ a σ : ℝ, 0 < Δμ →
    G Δμ a σ (rStar Δμ a σ) = (1 / 2) * Δμ * (rStar Δμ a σ) ^ 3)
\end{lstlisting}

\subsection{\texttt{P2.Afforestation}: cubic efficiency and seeding alignment}
Definitions: $\mathrm{eff}(a,\sigma,\Delta\mu,C)=C\,(3\Delta\mu/(2a\sigma))^3$ and $\eta(n,m)=\langle n,m\rangle^2/(\|n\|^2\|m\|^2)$.
\begin{proposal}
(i) For all $a,\sigma,\Delta\mu,t>0$ and all $C$: $\mathrm{eff}(a,\sigma,t\Delta\mu,C)=t^3\,\mathrm{eff}(a,\sigma,\Delta\mu,C)$.
(ii) In every real inner-product space $E$ (universe-polymorphic, as in the canon), for $n,m\ne0$: $0\le\eta(n,m)\le1$.
\end{proposal}
These restate the canon theorems
\texttt{establishment\_efficiency\_cubic\_in\_drive\_over\_tension} and \texttt{seeding\_efficiency\_eq\_cos2\_theta}
 (file sha256 \texttt{3ebe33bf\ldots 9526f}).
\begin{remark}
The second canon theorem is named ``$\mathrm{eff}=\cos^2\Theta$''. Its docstring also mentions $\eta=1$ iff $n\parallel m$, and $1-\eta=\sin^2\Theta$.
Its Lean conclusion proves only the bound $0\le\eta\le1$, and the proposal restates only that bound.
If the kernel uses the equality case, that case needs its own statement.
\end{remark}
Sanity examples pin $\mathrm{eff}=C/n^\star$ ($27/8$ at unit parameters), the factor $8$ when the drive is doubled, and the factor $8$ when $\sigma$ is halved.
They also pin $\eta=1$ for parallel and $\eta=0$ for orthogonal vectors in $\mathbb{R}^2$.
Negative controls: doubling the drive does not double the efficiency, and the unnormalised $\langle2,3\rangle^2=36$ exceeds $1$.
\begin{lstlisting}
def Statement : Prop :=
  (∀ a σ Δμ C t : ℝ, 0 < a → 0 < σ → 0 < Δμ → 0 < t →
    eff a σ (t * Δμ) C = t ^ 3 * eff a σ Δμ C) ∧
  (∀ (E : Type*) [NormedAddCommGroup E] [InnerProductSpace ℝ E] (n m : E),
    n ≠ 0 → m ≠ 0 → 0 ≤ eta n m ∧ eta n m ≤ 1)
\end{lstlisting}

\subsection{\texttt{P2.Harmonization}: one price suffices}
Definitions: $\mathrm{br}(\lambda)_i=b_i-\lambda/a_i$, $J(x)=\sum_i \tfrac{a_i}{2}(x_i-b_i)^2$,
$\lambda^\star=(\sum_i b_i-X)/\sum_i a_i^{-1}$ and $x^\star_i=\mathrm{br}(\lambda^\star)_i$.
\begin{proposal}
For all $N>0$, all $a,b\in\mathbb{R}^N$ with every $a_i>0$, and all $X$:
(i) $\sum_i \mathrm{br}(\lambda^\star)_i=X$; and
(ii) every $y$ with $\sum_i y_i=X$ satisfies $J(x^\star)\le J(y)$.
\end{proposal}
These restate \texttt{ght\_market\_clears} and \texttt{ght\_strong\_duality\_decentralized\_eq\_centralized}
(file sha256 \texttt{cfdead65\ldots 110aec}).
On the canon instance $N=2$, $a=b=(1,1)$, $X=0$, the sanity examples give $\lambda^\star=1$, $x^\star=(0,0)\ne b$, $J(x^\star)=1$, and $J(1,-1)=2$.
Negative controls:
\begin{itemize}
\item The infeasible ideal $y=b$ has $J=0<1$, so the feasibility hypothesis is essential.
\item The price $\lambda=0$ does not clear.
\item With $N=0$ the aggregate is $0\ne X=1$, so $N>0$ is essential.
\end{itemize}
\begin{lstlisting}
def Statement : Prop :=
  ∀ (N : ℕ) (a b : Fin N → ℝ), 0 < N → (∀ i, 0 < a i) → ∀ X : ℝ,
    (∑ i, br a b (lamStar a b X) i) = X ∧
    (∀ y : Fin N → ℝ, (∑ i, y i) = X → J a b (xStar a b X) ≤ J a b y)
\end{lstlisting}

\section{Tempo: a reconstructed statement (P2-T014)}\label{sec:tempo}
\subsection{What the kernel computes}
The registry entry for Tempo (\texttt{Tempo\_function.json}, viridis-canon commit \texttt{6a177d6b}) names six theorems in a
Lean module \texttt{StewardshipTempo}:
\texttt{Sigma\_strictConvexOn}, \texttt{tempo\_floor}, \texttt{tempo\_optimum}, \texttt{tempo\_unique\_minimizer},
\texttt{haste\_neglect\_cosh} and \texttt{governability\_tax}.
It lists no source file, and the only verification on record is a receipt from the retired Comparator.
The reference runner (\texttt{Tempo\_runner.py}) takes a haste cost $h>0$, a neglect cost $n>0$ and a cadence $t>0$, and returns
\[
t^\star=\sqrt{h/n},\qquad F=2\sqrt{hn},\qquad \Sigma(t)=\frac{h}{t}+nt,\qquad \frac{\Sigma(t)}{F},\qquad
\tau(t)=\cosh\log\frac{t}{t^\star},
\]
together with a floating-point flag that tests whether $t$ is close to $t^\star$.
These are the optimal cadence, the floor, the cost, the normalised cost and the haste--neglect tax.
The registry's published example is $(h,n,t)=(4,1,2)$, with expected outputs $t^\star=2$, $F=4$, $\Sigma=4$, normalised cost $1$ and tax $1$.

\subsection{Reconstruction and checks}
We first wrote one part per theorem name, following the task brief, and checked each part against the runner.
Write $s=\sqrt h$ and $r=\sqrt n$. Then $\sqrt{hn}=sr$ and $\sqrt{h/n}=s/r$, and
\[
\Sigma(t)-2sr=\frac{(s-rt)^2}{t}\qquad(t>0).
\]
This gives the floor, shows that it is attained at $t^\star=s/r$, and shows that it is attained nowhere else.
The identity $\cosh\log x=(x+x^{-1})/2$ for $x>0$ gives
$\cosh\log(t/t^\star)=(rt/s+s/(rt))/2=\Sigma(t)/(2sr)$. Since $\cosh\ge1$, with equality only at $0$, the tax is at least $1$ and equals $1$ only at $t^\star$.
Strict convexity of $h/t$ on $(0,\infty)$ follows, for $a+b=1$ with $a,b>0$ and $x\ne y$, from
$(ah/x+bh/y)(ax+by)=h+hab(x-y)^2/(xy)$.
The published example passes every check. All six parts as written in the brief are therefore true.

\paragraph{The deposited source.}
To cite the registry DOI, we resolved it. It points to a Zenodo record~\cite{HartTempo2026}, which lists a file \texttt{StewardshipTempo.lean}.
The record was modified on 14 November 2026, after the canon replay of 11 November that found no source.
We could not read the file directly in our sandbox. We obtained the signatures of its six theorems as a quotation relayed by
a web-fetching tool that summarises pages with a language model. We did not hash, compile or replay the file.
According to that quotation, in namespace \texttt{ViridisForge.Tempo} with cost $A/\tau+B\tau$ and hypotheses $A,B>0$:
\begin{itemize}
\item the first four theorems state exactly the brief's parts (1)--(4), with $A,B,\tau$ in place of $h,n,t$;
\item \texttt{haste\_neglect\_cosh} is the bare identity $(u+1/u)/2=\cosh\log u$ for $u>0$;
\item \texttt{governability\_tax} is $\Sigma(\tau)/(2\sqrt{AB})=\cosh\log(\sqrt{A/B}/\tau)$, with the ratio inverted relative to the runner;
\item no deposited theorem states that the tax is at least $1$ or characterises when it equals $1$.
\end{itemize}
The brief's part (5) is thus the deposited \texttt{governability\_tax}, up to the evenness of $\cosh$, and the brief's part (6) has no deposited counterpart.
Neither is wrong. They are attached to different names.

\subsection{The proposed statement}
We kept every part of the brief and added, under the matching name, each deposited signature that the brief did not already cover.
A proof of the sealed statement would then imply both the brief's six claims and the six deposited theorems.
\begin{proposal}[\texttt{P2.Tempo}]
For all real $h,n>0$:
(1) $\Sigma$ is strictly convex on $(0,\infty)$;
(2) $2\sqrt{hn}\le\Sigma(t)$ for every $t>0$;
(3) $\Sigma(\sqrt{h/n})=2\sqrt{hn}$;
(4) if $t>0$ and $\Sigma(t)=2\sqrt{hn}$ then $t=\sqrt{h/n}$;
(5) $(u+1/u)/2=\cosh\log u$ for every $u>0$, and $\Sigma(t)/(2\sqrt{hn})=\cosh\log(t/\sqrt{h/n})$ for every $t>0$;
(6) for every $t>0$, $\Sigma(t)/(2\sqrt{hn})=\cosh\log(\sqrt{h/n}/t)$, $\tau(t)\ge1$, and $\tau(t)=1\iff t=\sqrt{h/n}$.
\end{proposal}
\begin{lstlisting}
def Statement : Prop :=
  ∀ h n : ℝ, 0 < h → 0 < n →
    StrictConvexOn ℝ (Set.Ioi 0) (Sigma h n) ∧
    (∀ t : ℝ, 0 < t → floor h n ≤ Sigma h n t) ∧
    Sigma h n (optimum h n) = floor h n ∧
    (∀ t : ℝ, 0 < t → Sigma h n t = floor h n → t = optimum h n) ∧
    ((∀ u : ℝ, 0 < u → (u + 1 / u) / 2 = Real.cosh (Real.log u)) ∧
      (∀ t : ℝ, 0 < t → normalizedCost h n t = tax h n t)) ∧
    (∀ t : ℝ, 0 < t →
      normalizedCost h n t = Real.cosh (Real.log (optimum h n / t)) ∧
      1 ≤ tax h n t ∧ (tax h n t = 1 ↔ t = optimum h n))
\end{lstlisting}
Here \texttt{Sigma h n t = h / t + n * t}, \texttt{optimum h n = Real.sqrt (h / n)}, \texttt{floor h n = 2 * Real.sqrt (h * n)},
\texttt{normalizedCost = Sigma / floor} and \texttt{tax h n t = Real.cosh (Real.log (t / optimum h n))}.
All are defined inside \texttt{ViridisStatements.P2.Tempo}, and they unfold to the runner's formulas and to the deposited expressions.
If the brief's six-part form is preferred, deleting the two deposited conjuncts gives exactly that form.

\paragraph{Controls.} The module's examples are closed by \texttt{norm\_num} and short rewrites.
\texttt{decide} does not apply to $\mathbb{R}$.
\begin{itemize}
\item The published example: $t^\star=2$, $F=4$, $\Sigma(2)=4=F$, normalised cost $1$ and tax $1$.
\item A second witness: $(h,n)=(1,4)$ gives $t^\star=1/2$, $F=4$ and $\Sigma(1/2)=4$.
\item At $(4,1,1)$ the normalised cost is $5/4=\cosh\log 2$.
\end{itemize}
Negative controls:
\begin{itemize}
\item $(4,1,1)$ gives cost $5>4$.
\item $t=-2$ gives $\Sigma=-4<F$, so $t>0$ is needed.
\item $n=0$ breaks parts (5) and (6) under Lean's $x/0=0$: the normalised cost is $0$ while the tax is $1$, and the tax equals $1$ away from the ``optimum'' $0$.
\item $h=0$ makes $\Sigma$ linear, so $\Sigma(2)=(\Sigma(1)+\Sigma(3))/2$ and strict convexity fails.
\end{itemize}
The module compiles under the pinned toolchain. None of 17 one-shot tactics closes \texttt{Statement} after unfolding.
These include \texttt{simp}, \texttt{norm\_num}, \texttt{aesop}, \texttt{grind}, \texttt{nlinarith}, \texttt{exact?}, and a six-way split followed by \texttt{simp\_all}.
A private scratch derivation of the full statement compiled, which suggests the target is not false.
It is not offered as a proof. The proof belongs to the follow-up task P2-T015, against the sealed module (Section~\ref{sec:tempoproof}).
The statement was sealed unchanged, in the merged form above (decision D-021; sealed module sha256 \texttt{a966ef3b\ldots 2179}).

\section{Tempo: the Lean proof (P2-T015)}\label{sec:tempoproof}
\subsection{Result}
\begin{theorem}\label{thm:tempo}
The sealed statement \texttt{ViridisStatements.P2.Tempo.Statement}, parts (1)--(6) of the proposal in Section~\ref{sec:tempo}, holds.
We give a Lean proof in \texttt{proofs/P2.Tempo.lean}.
\end{theorem}
The sealed module (sha256 \texttt{a966ef3b1e3bb10cb9336686d71ec5c8b85c442fc1a21e6d3c2a13c7b6622179}) is identical to the
20 November proposal. The local checker \texttt{viridis-lean-check} reported \texttt{PASSED} on all five stages (denylist, compile,
statement match, axioms, kernel replay), with axioms \texttt{propext}, \texttt{Classical.choice} and \texttt{Quot.sound}.
The verifier awards the label independently.

\subsection{Proof structure}
The proof works in square-root coordinates. For $h,n>0$ put $s=\sqrt h$ and $r=\sqrt n$, so $h=s^2$ and $n=r^2$ with $s,r>0$.
A lemma \texttt{core} proves all six parts for $\Sigma(s^2,r^2,\cdot)$, and \texttt{main} rewrites $h$ and $n$ back.
This avoids reasoning about square roots of products and quotients inside the main argument; they are handled once, in two lemmas:
\[
\sqrt{s^2/r^2}=s/r,\qquad 2\sqrt{s^2r^2}=2sr .
\]
\begin{itemize}
\item \textbf{Parts (2)--(4)} follow from the gap identity $\Sigma(t)-2sr=(s-rt)^2/t$ (lemma \texttt{gap}, by \texttt{field\_simp} and \texttt{ring}).
The right side is nonnegative for $t>0$, giving the floor; it vanishes at $t=s/r$, giving attainment; and if it vanishes then
$s=rt$, giving uniqueness.
\item \textbf{Part (5)}: the identity $(u+1/u)/2=\cosh\log u$ is Mathlib's \texttt{Real.cosh\_log}. With it both
$\cosh\log(t/(s/r))$ and $\cosh\log((s/r)/t)$ reduce to $(rt/s+s/(rt))/2$, which equals $\Sigma(t)/(2sr)$ by field arithmetic
(lemma \texttt{norm\_eq\_cosh}). This gives both orientations at once, so the deposited form in part (6) needs no appeal to evenness of $\cosh$.
\item \textbf{Part (6)}: $1\le\cosh x$ is \texttt{Real.one\_le\_cosh}. For the equality case we did not use injectivity of $\cosh$ on $[0,\infty)$.
Instead, since the tax equals the normalised cost, tax $=1$ means $\Sigma(t)=2sr$, and part (4) gives $t=t^\star$; conversely part (3) gives tax $=1$ at $t^\star$.
\item \textbf{Part (1)}: Mathlib has convexity results for $x\mapsto x^{-1}$, but we proved strict convexity directly from the definition
of \texttt{StrictConvexOn}. With $b=1-a$, the linear term $nt$ contributes equally to both sides, so it suffices that
$h/(ax+by)<a\,h/x+b\,h/y$. Multiplying by $ax+by>0$ and using
$(ah/x+bh/y)(ax+by)=h+hab(x-y)^2/(xy)$, the extra term is strictly positive because $x\ne y$.
\end{itemize}
The helper lemmas are:
\begin{lstlisting}
/-- With `h = s ^ 2` and `n = r ^ 2`, the optimum is `s / r`. -/
lemma optimum_sq {s r : ℝ} (hs : 0 < s) (hr : 0 < r) :
    optimum (s ^ 2) (r ^ 2) = s / r := by
  rw [optimum, Real.sqrt_div' _ (by positivity), Real.sqrt_sq hs.le, Real.sqrt_sq hr.le]

/-- With `h = s ^ 2` and `n = r ^ 2`, the floor is `2 * s * r`. -/
lemma floor_sq {s r : ℝ} (hs : 0 < s) (hr : 0 < r) :
    floor (s ^ 2) (r ^ 2) = 2 * s * r := by
  rw [floor, ← mul_pow, Real.sqrt_sq (by positivity)]; ring

/-- Gap identity: `Sigma t - floor = (s - r t) ^ 2 / t`. -/
lemma gap {s r t : ℝ} (hs : 0 < s) (hr : 0 < r) (ht : 0 < t) :
    ViridisStatements.P2.Tempo.Sigma (s ^ 2) (r ^ 2) t - floor (s ^ 2) (r ^ 2) = (s - r * t) ^ 2 / t := by
  rw [floor_sq hs hr, ViridisStatements.P2.Tempo.Sigma]; field_simp; ring

/-- The cosh form of the normalised cost, in either orientation. -/
lemma norm_eq_cosh {s r t : ℝ} (hs : 0 < s) (hr : 0 < r) (ht : 0 < t) :
    normalizedCost (s ^ 2) (r ^ 2) t = Real.cosh (Real.log (t / (s / r))) ∧
    normalizedCost (s ^ 2) (r ^ 2) t = Real.cosh (Real.log (s / r / t)) := by
  rw [normalizedCost, floor_sq hs hr, ViridisStatements.P2.Tempo.Sigma,
    Real.cosh_log (by positivity), Real.cosh_log (by positivity)]
  constructor <;> (field_simp; try ring)
\end{lstlisting}
and the final step is:
\begin{lstlisting}
theorem main : ViridisStatements.P2.Tempo.Statement := by
  intro h n hh hn
  have eh : h = Real.sqrt h ^ 2 := (Real.sq_sqrt hh.le).symm
  have en : n = Real.sqrt n ^ 2 := (Real.sq_sqrt hn.le).symm
  have := core (Real.sqrt_pos.mpr hh) (Real.sqrt_pos.mpr hn)
  rw [← eh, ← en] at this
  exact this
\end{lstlisting}
The full file (127 lines) is in the reproducible package.

\subsection{Process}
The proof was written by hand from the derivation recorded on 20 November and passed the checker after two edits:
the name \texttt{Sigma} is ambiguous with Lean's built-in \texttt{\_root\_.Sigma} type once the statement namespace is opened, so it is
written fully qualified; and \texttt{field\_simp} alone closed one of the two \texttt{cosh} goals, so the trailing \texttt{ring} is wrapped in \texttt{try}.
Aristotle was not used.

\section{Results}
\begin{table}[h]
\centering\small
\begin{tabular}{@{}llll@{}}
\toprule
Kernel & Sealed statement & Session & Status of our candidate \\
\midrule
Carbon Continuity & \texttt{P2.CarbonContinuity} & 2026-10-29 & Lean proof, local check passed \\
Restoration & \texttt{P2.Restoration} & 2026-10-31 & Lean proof, local check passed \\
Afforestation & \texttt{P2.Afforestation} & 2026-11-01 & Lean proof, local check passed \\
Harmonization & \texttt{P2.Harmonization} & 2026-11-02 & Lean proof, local check passed \\
Tempo & \texttt{P2.Tempo} & 2026-11-21 & Lean proof, local check passed \\
\bottomrule
\end{tabular}
\caption{State of the P2 program. ``Local check passed'' means \texttt{viridis-lean-check} accepted the file, with axioms
\texttt{propext}, \texttt{Classical.choice} and \texttt{Quot.sound}. The verifier awards the labels.}
\end{table}
For Tempo, the result of 21 November is Theorem~\ref{thm:tempo}: the proof of the sealed \texttt{P2.Tempo} passed the local checker
after two small edits (Section~\ref{sec:tempoproof}). Aristotle was not used.

For Tempo, the results of 20 November were as follows.
\begin{itemize}
\item The proposed module compiles, and all its examples elaborate.
\item No probe closed \texttt{Statement}.
\item All six parts of the brief turned out to be true.
\item According to the relayed quotation, parts (1)--(4) agree exactly with the deposited source, and parts (5) and (6) were merged with it as described in Section~\ref{sec:tempo}.
\item Aristotle was not used.
\end{itemize}

For Harmonization, the results of 2 November are as follows.
The restructured proof of \texttt{P2.Harmonization} passed the local checker on the first attempt.
The verbatim canon proof also passed it. Part (ii) does not need $N>0$; a zero curvature breaks it.
Aristotle was not used.

For Afforestation, the results of 1 November were as follows.
The restructured proof of \texttt{P2.Afforestation} passed the local checker on the first attempt.
The verbatim canon proof also passed it. Neither part needs any of its hypotheses.
Aristotle was not used.

For Restoration, the results of 31 October were as follows.
The restructured proof of \texttt{P2.Restoration} passed the local checker on the first attempt.
The verbatim canon proof also passed it.
The half-barrier identity holds without its hypothesis $0<\Delta\mu$.
Aristotle was not used.

The 30 October proposals are summarised below.
\begin{table}[h]
\centering\small
\begin{tabular}{@{}p{0.17\textwidth}p{0.42\textwidth}lc@{}}
\toprule
Proposal & Canon theorems restated & Canon file sha256 (prefix) & Probes closed \\
\midrule
\texttt{P2.Restoration} & \texttt{critical\_nucleus}, \texttt{barrier\_half\_identity} & \texttt{667a826a} & 0/28 \\
\texttt{P2.Afforestation} & \texttt{establishment\_efficiency\_\ldots}, \texttt{seeding\_efficiency\_\ldots} & \texttt{3ebe33bf} & 0/28 \\
\texttt{P2.Harmonization} & \texttt{ght\_market\_clears}, \texttt{ght\_strong\_duality\_\ldots} & \texttt{cfdead65} & 0/28 \\
\texttt{P2.Tempo} & six registry names (Section~\ref{sec:tempo}); reconstructed & none in canon & 0/17 \\
\bottomrule
\end{tabular}
\caption{Proposed sealed statements. The first three are from 30 October, and \texttt{P2.Tempo} is from 20 November. All four modules compile under the pinned toolchain, and their sanity examples elaborate.
The probe column counts how many of our 28 one-shot tactics closed \texttt{Statement} locally. The verifier runs its own probes.}
\end{table}
We observe the following.
\begin{itemize}
\item For each proposal, a private port of the canon proof elaborated against a local build of the module.
It depends only on \texttt{propext}, \texttt{Classical.choice} and \texttt{Quot.sound}.
We do not present these as proofs of the proposals. The proofs belong to later tasks, against the sealed modules.
\item Several canon hypotheses are not needed mathematically but are kept for faithfulness:
\begin{itemize}
\item In Lean's convention $x/0=0$, the half-barrier identity also holds at $\Delta\mu=0$.
\item The cubic law is a ring identity, so it needs none of its positivity hypotheses.
\item The alignment bound holds trivially when $n=0$ or $m=0$, since then $\eta=0$.
\end{itemize}
\item One canon name over-promises: the ``$\cos^2\Theta$'' theorem proves only the bound $0\le\eta\le1$.
\end{itemize}

\section{What the statements settle and what they do not}
\begin{itemize}
\item \textbf{Settled, once sealed and proved:} the internal mathematics of each kernel's rule inside its stated model.
\item \textbf{Not settled:} any parameter value. This covers $\Delta\mu$, $\sigma$, $a$, $C$, the curvatures $a_i$, the ideals $b_i$ and the target $X$.
These come from field evidence and remain ASSUMED.
\item \textbf{Not settled:} whether the models are adequate. The models in question are classical nucleation for tree-cluster establishment and
quadratic dissipation for stewardship coordination. Nucleation rates, ecological temperature and stochastic establishment are outside the statements.
\item \textbf{Not settled:} code conformance. The kernels' runtime identifiers and TypeScript sources were not among this session's inputs.
\end{itemize}

\paragraph{Restoration in particular.} Theorem~\ref{thm:restoration} settles the following:
\begin{itemize}
\item Inside the cubic model with positive $\Delta\mu,a,\sigma$ there is exactly one barrier top on $r\ge0$, and it sits at
$r^\star=2a\sigma/(3\Delta\mu)$.
\item $G$ is at most $G(r^\star)$ everywhere on $r\ge 0$.
\item The barrier height is $\tfrac12\Delta\mu (r^\star)^3$.
\end{itemize}
So a kernel that computes $r^\star$ (or $n^\star=(r^\star)^3$) by this formula and calls ``go'' above it and ``no-go'' below it
is applying a correct consequence of the model.
The theorem does \emph{not} settle any of the following, all of which stay ASSUMED:
\begin{itemize}
\item the field values of the driving force $\Delta\mu$, the edge ``surface tension'' $\sigma$ and the shape factor $a$, and
their uncertainty. Because $n^\star$ scales as $(a\sigma/\Delta\mu)^3$, relative errors in these inputs are roughly tripled in $n^\star$;
\item whether classical nucleation theory, with an $n^{2/3}$ edge cost, describes tree-cluster establishment at a given site;
\item the time-scale and probability layer: the Arrhenius nucleation rate, the ecological temperature, the stochastic committor
that replaces the sharp threshold in practice, and any planting deadline;
\item the seed-allocation result \texttt{broadcast\_suboptimal}, which is in the same canon file but not in the sealed statement;
\item the translation from $r$ to cluster size $n$ (a short argument, not formalised here);
\item that the deployed code computes $r^\star$ with this formula and applies the stated comparison.
\end{itemize}

\paragraph{Afforestation in particular.} Theorem~\ref{thm:afforestation} settles that, inside the model,
doubling the ecological drive multiplies seeding efficiency by exactly $8$, and more generally by $t^3$; so a kernel that ranks
plans by $\mathrm{eff}$ will see drive (and, by the same algebra, edge hostility) changes amplified cubically. It also settles that
the alignment score is a number in $[0,1]$, so the forcing debt $1-\eta$ is also in $[0,1]$.
It does \emph{not} settle, and the following stay ASSUMED:
\begin{itemize}
\item the field values of $a$, $\sigma$, $\Delta\mu$ and the canopy yield $C$; because of the cube, relative errors in
$\Delta\mu/(a\sigma)$ are roughly tripled in $\mathrm{eff}$;
\item whether ``canopy per seed'' $C/n^\star$ is the right ranking criterion, and whether the nucleation model applies to the site;
\item the equality case $\eta=1\iff n\parallel m$ and the scale invariance of $\eta$;
\item how the optimal seeding vector $m$ is computed, which is outside the statement;
\item the other canon results in the same file (the homogeneous-limit barrier identity and the information-rate ceiling);
\item that the deployed code computes $\mathrm{eff}$ and $\eta$ by these formulas.
\end{itemize}

\paragraph{Harmonization in particular.} Theorem~\ref{thm:harmonization} settles that, inside the separable quadratic model
with positive curvatures, broadcasting the single number $\lambda^\star$ is enough: if every steward best-responds to it, the
shared target is met exactly, and no central assignment meeting the target dissipates less. A kernel that computes $\lambda^\star$
by this formula and publishes it is therefore applying a correct consequence of the model.
It does \emph{not} settle, and the following stay ASSUMED:
\begin{itemize}
\item the field values of the curvatures $a_i$, the ideal efforts $b_i$ and the target $X$; errors in $a_i$ enter $\lambda^\star$
through $\sum_i 1/a_i$, so a single steward with badly underestimated curvature can dominate the price;
\item whether stewards' costs are quadratic and separable, and whether they actually best-respond to a price;
\item any dynamics: the canon's primal--dual (t\^atonnement) convergence result is deferred in the canon itself and is not sealed,
so a kernel that discovers the price iteratively rather than by the closed form has no formal backing for convergence here;
\item uniqueness of the optimum and the other canon results in the same file;
\item that the deployed code computes $\lambda^\star$ and $\mathrm{br}$ by these formulas.
\end{itemize}

\paragraph{Tempo in particular.} Theorem~\ref{thm:tempo} settles the following inside the model
$\Sigma(t)=h/t+nt$ with constant $h,n>0$:
\begin{itemize}
\item the runner's optimum, floor and tax are mutually consistent;
\item $t^\star=\sqrt{h/n}$ is the unique cost-minimising cadence;
\item the normalised cost of any other cadence equals $\cosh\log(t/t^\star)>1$.
\end{itemize}
Because $\cosh$ is even, being a factor $k$ too fast costs exactly as much as being a factor $k$ too slow.
It does \emph{not} settle, and the following stay ASSUMED:
\begin{itemize}
\item \textbf{The inputs are declared, not established.} The cost model, and its two coefficients $h$ and $n$, are what a user
declares when calling the kernel. The theorem says what follows if they are right; it says nothing about whether they are.
\item \textbf{Stewardship outcomes.} Nothing here shows that following $t^\star$ improves any ecological or programme outcome.
\item \textbf{The retired receipt.} The proof replaces, rather than confirms, the Comparator receipt: we do not know what that
receipt checked, because the module it refers to is not on file in the canon repository.
\item \textbf{The model itself.} The form $h/t+nt$ with constant haste and neglect costs is an assumption about a
stewardship programme, not a result. Nothing here shows that real costs scale like $1/t$ and $t$, or that $h$ and $n$ are
measurable or stable over time.
\item \textbf{The units of the inputs.} The registry gives both $h$ and $n$ the unit ``cost per cycle''. For $h/t+nt$ to add like
terms, $n$ must instead carry cost per time squared when $h$ is a cost.
\item \textbf{The kernel's floating-point behaviour.} This covers the \texttt{at\_optimum} tolerance and rounding to 12 digits.
\item \textbf{Code conformance.} The statement does not cover whether the deployed TypeScript runtime matches the runner.
\end{itemize}

\section{Limitations}
The mathematics is elementary, and its value lies in being faithful and re-checkable.
Faithfulness was judged by hand against the canon text. No tool checks that a proposal matches its canon theorem.
Our probe list is our own guess at the verifier's and may differ from it.
\texttt{P2.Afforestation} is universe-polymorphic, following the canon's \texttt{Type*}. We checked locally that a polymorphic
candidate \texttt{theorem main : Statement} and the verifier's statement-match template both elaborate.
A universe-monomorphic seal (\texttt{Type}) would be a mild weakening.
The project README asks readers to consult a guidance file (\texttt{STATEMENT\_PROPOSAL.md}) that was not among the inputs.
We followed the brief and the format of the existing sealed module instead. That file was again absent on 31 October, 1 November, 2 November and 20 November.
For Tempo, the statement is a reconstruction. Its comparison with the deposited \texttt{StewardshipTempo.lean} rests on a
quotation of the theorem signatures relayed by a summarising web tool, not on a byte-level read. That file's hash, its
compilation and its axioms are unknown to us, and the Zenodo record notes that a newer version exists.
For Restoration, the sealed statement quantifies over the cube-root variable $r$, while field practice counts trees $n$.
The step between the two is informal in this paper.

\section{Reproducibility}
Toolchain \texttt{leanprover/lean4:v4.28.0}, Mathlib commit \texttt{8f9d9cff6bd728b17a24e163c9402775d9e6a365}.
Canon repository commit \texttt{80d32904d751b45f415e23f872510cc2eddca3ff}. The canon file hashes are listed in the proposals' rationale files.
Each proposal (in the outputs of run 2026-10-30) compiles with \texttt{lean proposals/P2.}\textit{Name}\texttt{.lean}.
The sealed Carbon Continuity statement has sha256 \texttt{a5a6e8c6\ldots ffded}. To check its proof, run
\texttt{viridis-lean-check proofs/P2.CarbonContinuity.lean P2.CarbonContinuity}.
The sealed Restoration statement has sha256 \texttt{fd40b113c06beaed7940b2bf84d931a8b9179dcabaf3c2d27f7300dc27ccff18},
and the canon file it restates has sha256 \texttt{667a826a2e5befac67964b69670435362de2d3d792e163807c3bb7c2aec28132}.
To check the proof, run \texttt{viridis-lean-check proofs/P2.Restoration.lean P2.Restoration}.
The sealed Afforestation statement has sha256 \texttt{2fc1cbb05d7780f266f9da9c00941ca1b7d8ebc703a11a25fc9bb23237a59436},
and the canon file it restates has sha256 \texttt{3ebe33bfd6bd74124e87b5d7e73448db58e1459207fd6ff5c00a90fc3ac9526f}.
To check the proof, run \texttt{viridis-lean-check proofs/P2.Afforestation.lean P2.Afforestation}.
The sealed Harmonization statement has sha256 \texttt{48934ecb301f858b8332eea15194ff42a119fd7c1c7a7402089467a4994b3dd7},
and the canon file it restates has sha256 \texttt{cfdead658b1373cd311ac816f88867c5fd429a708e813a23a3293c901d110aec}.
To check the proof, run \texttt{viridis-lean-check proofs/P2.Harmonization.lean P2.Harmonization}.
The Tempo proposal compiles with \texttt{lean proposals/P2.Tempo.lean} (outputs of run 2026-11-20). Its context files
\texttt{Tempo\_function.json} and \texttt{Tempo\_runner.py} have sha256 prefixes \texttt{3a8c5312} and \texttt{a98b5c65}, which match
the project manifest. The deposited source was consulted only through the Zenodo record~\cite{HartTempo2026}.
The Tempo statement was sealed with sha256 \texttt{a966ef3b1e3bb10cb9336686d71ec5c8b85c442fc1a21e6d3c2a13c7b6622179}.
To check its proof, run \texttt{viridis-lean-check proofs/P2.Tempo.lean P2.Tempo}; the check takes well under a minute.
No automated prover jobs (Aristotle) and no external data were used in any of the seven sessions.

\section{Changelog}
\begin{itemize}
\item 2026-10-29 (P2-T001): First version. Ported the canon threshold proof to the sealed statement
\texttt{P2.CarbonContinuity}, and wrote the scope statement.
\item 2026-10-30 (P2-T002): Proposed sealed statements \texttt{P2.Restoration}, \texttt{P2.Afforestation} and \texttt{P2.Harmonization},
with rationales, sanity examples and negative controls. Added a method section for writing statements.
Condensed the Carbon Continuity material into its own section, and retitled the paper to cover the P2 program.
\item 2026-10-31 (P2-T003): Gave a Lean proof of the sealed statement \texttt{P2.Restoration}, as a restructured port of the canon
proof. Added the Restoration section, covering the proof, its relation to the canon, an unused hypothesis, and the $r$-versus-$n$ remark.
Added a scope paragraph on what the theorem does not settle for the Restoration kernel. Added a status table, and updated the title and abstract.
\item 2026-11-01 (P2-T004): Gave a Lean proof of the sealed statement \texttt{P2.Afforestation}, porting the canon proof with
an explicit squaring step in place of \texttt{nlinarith}. Added the Afforestation section (proof, relation to canon, unused hypotheses,
the over-promising name), a scope paragraph for the Afforestation kernel, and updated the status table, title and abstract.
\item 2026-11-02 (P2-T005): Gave a Lean proof of the sealed statement \texttt{P2.Harmonization}, porting the canon proof with an
explicit linearity lemma and explicit cancellation in place of \texttt{simp +decide}. Added the Harmonization section (proof,
relation to canon, the role of $N>0$ and of positive curvature), a scope paragraph for the Harmonization kernel, and updated the
status table, title, abstract and reproducibility section.
\item 2026-11-20 (P2-T014): Proposed the sealed statement \texttt{P2.Tempo} for the Tempo cadence kernel, reconstructed from the
reference runner. Compared it with the theorem signatures of a deposited \texttt{StewardshipTempo.lean} found on Zenodo, and
merged the deposited forms of \texttt{haste\_neglect\_cosh} and \texttt{governability\_tax} into parts (5) and (6).
Added the Tempo section, a Tempo scope paragraph, a limitations note on the deposit's provenance, and a bibliography entry.
Retitled the paper and updated the abstract, introduction, background and tables.
\item 2026-11-21 (P2-T015): Gave a Lean proof of the sealed statement \texttt{P2.Tempo}. Added the Tempo proof section, the Tempo
row of the status table, the retired-receipt and declared-input points to the Tempo scope paragraph, and the reproducibility entry.
Retitled the paper as the Tempo re-certification note and rewrote the abstract.
\end{itemize}

\bibliographystyle{plainnat}
\bibliography{refs}
\end{document}
